\subsection{Finite-prevalence comparisons with the limiting risks}
\label{sec:oracle-finite}

We compare \cref{thm:oracle-limits,cor:oracle-allocation} with exact moments
under the same population model. The saved study contains 43 signal, noise,
and target-budget settings at six prevalences from $0.04$ to $2\times10^{-6}$.
The settings include both finite-prevalence theorem anchors at $\tau=1,2,4$;
common signals at 90\%, 100\%, and 110\% of the predicted crossing; eight
allocations with $s/4+r=1$ and $\tau=2$; four noise levels; points around U's
rare-learning threshold; budgets from $0.01$ to $1000$; and zero-signal and
zero-noise endpoints. These are selected slices rather than uniform coverage
of the surrounding parameter ranges.
\Cref{fig:oracle-finite} summarizes the comparisons.

Each scaling minimizes expected rare risk over separate finite grids. The fine
rate grid contains 2049 fractions of its mean-square stability limit, combining
1025 linear and 1025 logarithmic points, with zero and positive fractions from
$10^{-10}$ to $1-10^{-6}$. Limiting optimal rates are not inserted into the
grid, and rates outside the stable interval are not searched. Stopping grids
contain initialization, every integer through 64, all caps, and rounded
geometric checkpoints up to the largest cap. They have 663 to 883 distinct
stops, depending on prevalence. The largest cap is $500{,}001{,}001$.
Matrix powers of the affine recurrence evaluate every candidate without
Monte Carlo noise.

Caps round $\tau/b$ upward, except at target $\tau=4$, where they round
downward to stay within the finite-prevalence theorem's range. Primary errors
use $L_A(Mb)$, and the saved results also retain predictions at the nominal
target. At $q=0.04$, target $\tau=0.01$ gives one observation and actual
$Mb=0.0384$, so the coarse-budget discrepancy includes appreciable rounding.

Coarse grids take every second generating rate and geometric stop, retaining
the early integer stops and all caps. Separate rate refinement, stop refinement,
and joint refinement assess search resolution. The largest decrease under joint
refinement is $1.362\times10^{-5}$. Stop refinement alone changes no selected
risk: every positive-stop selection uses its cap, and the others use
initialization. No positive-stop selection reaches a rate-grid edge. These
observations do not bound the error relative to the unrestricted oracle or
prove that full-cap training is always optimal at finite prevalence.

\begin{figure}[t]
  \centering
  \includegraphics{../figures/oracle-finite/limit-comparison.pdf}
  \caption{
    Exact finite-grid selections approach the limiting oracle predictions.
    (a) Largest absolute risk discrepancy from $L_A(Mb)$ across the 43 settings
    at each prevalence. (b) Allocation at fixed initial limiting subgroup risk.
    (c) Noise can change the preferred scaling twice, with a finite-prevalence
    disagreement at high noise. (d) The budget comparison approaches the ordered
    large-budget risk ratio $4/5$. Colored curves connect finite-grid results;
    black dashed curves connect limiting predictions at the sampled settings.
    Panels (b)--(d) use nominal target budgets, so coarse caps also reflect
    rounding. Lines do not certify intervening crossing locations.
  }
  \label{fig:oracle-finite}
\end{figure}

Both theorem orderings hold in every sampled anchor comparison. At
$q=2\times10^{-6}$, the largest anchor errors are $2.987\times10^{-6}$ for U
and $1.536\times10^{-6}$ for S. Across all 43 settings, the corresponding
maxima are $2.01631\times10^{-4}$ and $2.27583\times10^{-6}$; the largest U
error occurs at target $\tau=0.01$. Every setting with a nonzero nominal
limiting gap has the predicted sign at the smallest prevalence. The sampled
common-signal crossings are bracketed by 90\% and 110\% of their limiting
values, and the allocation slice recovers the ordering on both sides of
$\theta_*\approx0.83770278$. Cases exactly at limiting ties remain sensitive
to rounding and search resolution and do not locate a finite-prevalence boundary.

The U endpoint has a concrete finite-prevalence interpretation. At
$\tau=2,s=1,r=0.01,\nu=1$, the limiting rate is zero and $L_U=r/2=0.005$.
The selected finite rate nevertheless remains positive, falling from $0.07003$
at $q=0.02$ to $2.866\times10^{-5}$ at $q=2\times10^{-6}$, while the risk
falls from $0.01570$ to $0.00500127$. These rates learn the common signal;
literal zero-rate training would retain its error.

\paragraph{A finite-prevalence disagreement.}

For $s=0.1,r=1$, target $\tau=2$, and $\nu=100$, the limiting risks favor U:
$L_U=0.500000$ and $L_S\approx0.502306$. The finite-grid SD-minus-U gaps are
instead $-0.00568756$ at $q=0.02$, $-0.00563282$ at $q=0.002$, and
$-0.000608220$ at $q=0.0002$. They become positive at $q=0.00002$ and
$q=0.000002$, with gaps $0.00180198$ and $0.00223702$. These discrepancies
exceed the observed refinement changes, showing that the limiting decision
can be inaccurate at finite prevalence. They do not certify a reversal of
the unrestricted finite-prevalence oracles.

At $s=r=\nu=1$ and the smallest prevalence, U wins at target $\tau=1$, while
S wins at $\tau=10,100,1000$. At $\tau=1000$, the finite risk ratio is
$0.839843$ and the limiting ratio is $0.839507$, both still above $4/5$.
These comparisons preserve the distinction between the prevalence limit and
the subsequent large-budget limit.

\paragraph{Reproduction and independent replay.}

The settings, grids, all four searches, selected risks, rates, stops, and
summaries are saved under \path{results/oracle-finite/}. The checker
\path{checks/oracle_finite.py} independently constructs observation-level
tensor updates and tests all four nested searches on a small problem. It also
replays all 2064 saved selections with full tensor updates and long-double
matrix powers; the recorded maximum discrepancy is $7.38\times10^{-11}$.
Run the following commands to regenerate the searches and summaries, reconstruct
every selection, and plot the saved results, respectively:
\begin{verbatim}
task experiment:oracle-finite
task check:oracle-finite -- --rerun
task figure:oracle-finite
\end{verbatim}
Without \texttt{--rerun}, the checker replays saved selections without repeating
the full searches. Summarizing, plotting, and building the paper never launch
these experiments.

\subsection{Introductory signal-allocation illustration}
\label{sec:signal-allocation}

\Cref{fig:intro-signal-crossing} uses a separate finite-grid study with the same
fresh-sample model, $q = 10^{-4}$, population preprocessing, zero initialization,
and fixed intercept. We set $(\beta_C^*)^2 = 1.01(1 - f)$ and
$(\beta_R^*)^2 = 1.01f$. The study holds the squared coefficient sum fixed;
\cref{cor:oracle-allocation} instead holds initial limiting subgroup risk fixed.
The grid contains 41 equally spaced values of $f$ in $[0,1]$, plus $1/101$ and
$100/101$, which recover the two signal pairs in
\cref{fig:conference-tuned-intro}. We cross these 43 allocations with noise
variances $1/4$, $1$, and $4$, and 21 integer caps evenly spaced up to rounding
from $10{,}002$ to $40{,}004$. The introductory panel displays noise variance
one and cap $M = 20{,}503$, so $Mq = 2.0503$. Its middle tick corresponds to
$f = 1/2$, where the two coefficient magnitudes are equal.

For each scaling, the fine rate grid contains zero and 1025 geometric values of
$\alpha$ from $10^{-6}$ to $2$. Stopping candidates contain zero, every cap, and
the distinct rounded values of 1025 geometric points from one to the largest
cap. Powers of the affine diagonal-moment recurrence give exact risks at each
candidate. The coarse search keeps every second positive rate and every second
geometric stopping point, together with zero and all caps. Joint refinement
changes any reported risk by less than $3 \times 10^{-5}$. All positive-stop
selections have interior rates. Replaying 112 selected endpoints with the
separately enumerated full tensor operator gives maximum absolute risk
discrepancy below $10^{-9}$. These checks assess the finite-grid calculation;
they do not certify unrestricted oracle risks or the crossing location.

The file \path{results/signal-allocation/summary.json} retains the grids, selected
rates and stops, risks, and checks. Run \texttt{task} \path{experiment:signal-allocation}
to regenerate the study, \texttt{python -m experiments.signal\_allocation.run
--check} to check the saved results, and \texttt{task} \path{figure:intro} to
plot them without rerunning the experiment.
